% Part: normal-modal-logic
% Chapter: frame-definability
% Section: frames

\documentclass[../../../include/open-logic-section]{subfiles}

\begin{document}

\olfileid{nml}{frd}{fra}

\olsection{ફ્રેમો}

\begin{defn}
  \emph{ફ્રેમ} એ જોડ $\mModel{F} = \tuple{W,R}$ છે, જ્યાં
  $W$ જગતોનો અખાલી ગણ અને $R$ એ~$W$ પરનો
  દ્વિસ્થાની સંબંધ છે. નિદર્શન $\mModel{M}$ ફ્રેમ
  $\mModel{F} = \tuple{W,R}$ \emph{પર આધારિત} છે
  ત્યારે અને માત્ર ત્યારે જ જ્યારે કોઈ નિમણૂક $V$
  માટે $\mModel{M} = \tuple{W, R, V}$ હોય.
\end{defn}

\begin{defn}
  $\mModel{F}$ ફ્રેમ હોય, તો $!A$ એ
  \emph{$\mModel{F}$માં માન્ય છે}, સંકેતમાં
  $\mModel{F} \Entails !A$, એમ કહીએ જો
  ફ્રેમ~$\mModel{F}$ પર આધારિત દરેક નિદર્શન~$\mModel{M}$
  માટે $\mSat{M}{!A}$ હોય.
  
  $\mClass{F}$ ફ્રેમોનો વર્ગ હોય, તો $!A$ એ
  \emph{$\mClass{F}$માં માન્ય છે}, સંકેતમાં
  $\mClass{F} \Entails !A$, ત્યારે અને માત્ર ત્યારે જ
  જ્યારે દરેક ફ્રેમ~$\mModel{F} \in \mClass{F}$ માટે
  $\mModel{F} \Entails !A$ હોય.
\end{defn}

ફ્રેમો રસપ્રદ હોવાનું કારણ એ છે કે સૂત્ર-ઢાંચા અને
સુલભતા સંબંધ~$R$ના ગુણધર્મો વચ્ચેની અનુરૂપતા
ફ્રેમોના સ્તરે છે, \emph{નિદર્શનોના સ્તરે નથી}.
ઉદાહરણ તરીકે, \Ax{T} બધાં સ્વવાચક નિદર્શનોમાં
સત્ય હોવા છતાં જેમાં \Ax{T} સત્ય હોય તે દરેક
નિદર્શન સ્વવાચક નથી. પરંતુ \Ax{T} બધાં સ્વવાચક
\emph{ફ્રેમો}માં \emph{માન્ય} છે અને, વધુમાં, જેમાં
\Ax{T} માન્ય હોય તે દરેક ફ્રેમ સ્વવાચક છે.

\begin{rem}
ફ્રેમોના વર્ગમાં માન્યતા એ નિદર્શનોના વર્ગમાં
માન્યતાનો વિશેષ કિસ્સો છે: $\mClass{F} \Entails !A$
ત્યારે અને માત્ર ત્યારે જ જ્યારે
$\mClass{C} \Entails !A$, જ્યાં $\mClass{C}$
એ~$\mClass{F}$ની કોઈ ફ્રેમ પર આધારિત બધા
નિદર્શનોનો વર્ગ છે.

સ્પષ્ટ છે કે કોઈ !!a{formula} અથવા સૂત્ર-ઢાંચો
માન્ય હોય, એટલે કે \emph{બધાં} નિદર્શનોના વર્ગ
સાપેક્ષે માન્ય હોય, તો તે ફ્રેમોના કોઈપણ
વર્ગ~$\mClass{F}$ સાપેક્ષે પણ માન્ય છે.
\end{rem}

\end{document}
